% ============================================================
% SECTION 5 — Breaking Down the Barriers (Ch. 12)
% ============================================================
\section{Breaking Down the Barriers}

\begin{frame}{Love Is What Breaks You Out of Yourself}
\bigquote{Self-love is a limitation and a defensive barrier; love completely breaks down this defensive barrier.}{Mutahhari}
{\footnotesize Timid. Cut off. Alone. --- until love breaks the wall.}
\bigquote{The struggle with self-love is the struggle with the limitations of the self.}{Mutahhari}
\begin{itemize}[itemsep=3pt]
  \item Not less self --- a \emph{bigger} self
  \item ``The personality must expand to take in every other human being.''
\end{itemize}
\bigquote{The reform of man does not lie in reducing him, it lies in perfecting and adding to him.}{Mutahhari}
\delivernote{Clarify live: this isn't about hating yourself -- it's expansion, not erasure.}

\end{frame}

% ---------------- Q&A ----------------
\qaframe{Breaking Down the Barriers}{%
  \qa{What does Mutahhari mean by ``self-love'' as a ``defensive barrier''?}
     {As long as a person is locked inside self-love, they remain weak -- timid, avaricious, quick-tempered, cut off from real vitality. Love breaks that barrier down, connecting the person to what is outside themselves.}
  \qa{Does Mutahhari say the goal is to eliminate self-regard? What is the real goal instead?}
     {No -- eliminating legitimate self-regard (``amour-propre'') is not even coherent. The real goal is to \emph{expand} the self so its boundary takes in other human beings, and ideally the whole of creation, rather than treating everything outside ``me'' as foreign.}
  \qa{In his own words, what is the difference between ``reducing'' a person and ``perfecting'' a person?}
     {Reform does not mean stripping away parts of a person's existence -- it lies in perfecting and adding to him, in growth and expansion, not in decrease or reduction.}
}
